\SourcePage{36}{41}

\section{Photon polarization}

For a photon, the polarization vector $\mathbf{e}$ serves as the ``spin part'' of the wave function (subject to the qualifications concerning the notion of photon spin stated in \S~6).

The various possibilities for photon polarization are the same as the possible polarization types of a classical electromagnetic wave (see II, \S~48).

An arbitrary polarization $\mathbf{e}$ can be represented as a superposition of two mutually orthogonal polarizations $\mathbf{e}^{(1)}$ and $\mathbf{e}^{(2)}$, chosen in some definite manner ($\mathbf{e}^{(1)}\mathbf{e}^{(2)*}=0$). In the expansion
\begin{equation}
\mathbf{e}=e_1\mathbf{e}^{(1)}+e_2\mathbf{e}^{(2)}
\tag{8.1}
\end{equation}
the squared moduli of $e_1$ and $e_2$ give the probabilities that the photon has polarization $\mathbf{e}^{(1)}$ or $\mathbf{e}^{(2)}$.

The two chosen polarizations may be taken as mutually perpendicular linear polarizations. Alternatively, an arbitrary polarization can be resolved into two circular polarizations rotating in opposite senses. We denote the right- and left-circular polarization vectors by $\mathbf{e}^{(+1)}$ and $\mathbf{e}^{(-1)}$, respectively; in a coordinate system $\xi\eta\zeta$ whose $\zeta$ axis lies along the photon direction $\mathbf{n}=\mathbf{k}/\omega$,
\begin{equation}
\mathbf{e}^{(+1)}=-\frac{i}{\sqrt{2}}(\mathbf{e}^{(\xi)}+i\mathbf{e}^{(\eta)}),
\qquad
\mathbf{e}^{(-1)}=\frac{i}{\sqrt{2}}(\mathbf{e}^{(\xi)}-i\mathbf{e}^{(\eta)}).
\tag{8.2}
\end{equation}

The availability of two distinct photon polarizations (at a specified momentum) means, equivalently, that every momentum eigenvalue is doubly degenerate. This fact is closely connected with the vanishing of the photon mass.

A freely moving particle of nonzero mass always has a rest frame. Evidently, it is in that frame that the particle's own symmetry properties become manifest. One must then consider symmetry under every possible rotation about the center\SourcePage{37}{42} (that is, under the entire spherical-symmetry group). The particle's symmetry with respect to this group is characterized by its spin $s$, which fixes the degeneracy multiplicity (the number $2s+1$ of distinct wave functions transformed into one another). In particular, a particle with a vector wave function (three components) has spin~1.

A zero-mass particle, by contrast, has no rest frame: it travels at the speed of light in every reference frame. Such a particle always singles out a direction in space, namely the direction of its momentum vector $\mathbf{k}$ (the $\zeta$ axis). Hence there is no symmetry under the full group of three-dimensional rotations; only axial symmetry about the distinguished axis remains.

With axial symmetry, the only conserved quantity is the particle's \emph{helicity}, the projection of its angular momentum on the $\zeta$ axis; denote it by $\lambda$.\footnote{This is to be distinguished from the projection $m$ of angular momentum on a specified direction in space (the $z$ axis), discussed in the preceding section.} If symmetry under reflections in planes containing the $\zeta$ axis is required as well, states differing by the sign of $\lambda$ are mutually degenerate; consequently, $\lambda\ne0$ gives a twofold degeneracy.\footnote{The electronic terms of a diatomic molecule are classified in the same way (see III, \S~78).} A photon state of specified momentum $\mathbf{k}$ belongs to one of these types of doubly degenerate states. Its ``spin'' wave function is a vector $\mathbf{e}$ in the $\xi\eta$ plane; the two components of this vector transform into each other under all rotations about the $\zeta$ axis and under reflections in planes containing that axis.

The different polarization cases correspond in a definite way to the possible helicity values. This correspondence follows from the formulas III, (57.9), which relate the components of a vector wave function to those of the equivalent second-rank spinor.\footnote{Recall that contravariant spinor components correspond to the components of a wave function regarded as probability amplitudes for the various values of the particle's angular-momentum projection, which are precisely the quantities at issue here.} Projections $\lambda=+1$ and $-1$ correspond to vectors $\mathbf{e}$ for which only $e_\xi-ie_\eta$ or $e_\xi+ie_\eta$, respectively, is nonzero; thus $\mathbf{e}=\mathbf{e}^{(+1)}$ or $\mathbf{e}=\mathbf{e}^{(-1)}$. In other words, $\lambda=+1$ and $-1$ correspond to right- and left-circular photon polarization (the same result will be obtained in \S~16 by directly calculating the eigenfunctions of the spin-projection operator). The photon angular momentum projected on its direction of motion can therefore take only the two values $\pm1$; zero cannot occur.

\SourcePage{38}{43}
The state of a photon with specified momentum and polarization is a pure state (in the sense explained in III, \S~14); it is described by a wave function and supplies a complete quantum-mechanical specification of the particle's state. ``Mixed'' photon states are possible as well; they correspond to a less complete description, given only by a density matrix rather than by a wave function.

Consider a photon state that is mixed with respect to polarization but has a definite momentum $\mathbf{k}$. In such a state (called a state of \emph{partial polarization}) a ``coordinate'' wave function exists.

The photon polarization density matrix is a second-rank tensor $\rho_{\alpha\beta}$ in the plane perpendicular to $\mathbf{n}$ (the $\xi\eta$ plane; the indices $\alpha$, $\beta$ take only two values). This tensor is Hermitian:
\begin{equation}
\rho_{\alpha\beta}=\rho^*_{\beta\alpha},
\tag{8.3}
\end{equation}
and is normalized by
\begin{equation}
\rho_{\alpha\alpha}=\rho_{11}+\rho_{22}=1.
\tag{8.4}
\end{equation}
By (8.3), the diagonal components $\rho_{11}$ and $\rho_{22}$ are real, and (8.4) determines either one from the other. The component $\rho_{12}$ is complex, while $\rho_{21}=\rho^*_{12}$. Thus the density matrix is specified by three real parameters in all.

Once the polarization density matrix is known, the probability that the photon has any particular polarization $\mathbf{e}$ can be found. It is given by the ``projection'' of the tensor $\rho_{\alpha\beta}$ on the direction of $\mathbf{e}$, namely
\begin{equation}
\rho_{\alpha\beta}e^*_\alpha e_\beta.
\tag{8.5}
\end{equation}
Thus $\rho_{11}$ and $\rho_{22}$ are the probabilities of linear polarization along the $\xi$ and $\eta$ axes. Projection onto the vectors (8.2) gives the probabilities of the two circular polarizations:
\begin{equation}
\frac{1}{2}[1\pm i(\rho_{12}-\rho_{21})].
\tag{8.6}
\end{equation}

Both in form and in content, the properties of $\rho_{\alpha\beta}$ agree with those of the tensor $J_{\alpha\beta}$ used for partially polarized light in the classical theory (see II, \S~50). We recall some of these properties here.

For a pure state of definite polarization $\mathbf{e}$, the tensor $\rho_{\alpha\beta}$ reduces to products of components of $\mathbf{e}$:
\begin{equation}
\rho_{\alpha\beta}=e_\alpha e^*_\beta.
\tag{8.7}
\end{equation}
Its determinant is then $|\rho_{\alpha\beta}|=0$. In the opposite case of an unpolarized photon, all polarization directions are equally probable, so that
\begin{equation}
\rho_{\alpha\beta}=\delta_{\alpha\beta}/2,
\tag{8.8}
\end{equation}
and $|\rho_{\alpha\beta}|=1/4$.

\SourcePage{39}{44}
In the general case, partial polarization is conveniently described by three real Stokes parameters $\xi_1$, $\xi_2$, $\xi_3$,\footnote{The notation for these parameters must not be confused with that for the $\xi$ axis.} in terms of which the density matrix is
\begin{equation}
\rho_{\alpha\beta}=\frac{1}{2}
\begin{pmatrix}
1+\xi_3 & \xi_1-i\xi_2\\
\xi_1+i\xi_2 & 1-\xi_3
\end{pmatrix}.
\tag{8.9}
\end{equation}
Each parameter ranges from $-1$ to $+1$. In the unpolarized state $\xi_1=\xi_2=\xi_3=0$; for a fully polarized photon, $\xi_1^2+\xi_2^2+\xi_3^2=1$.

The parameter $\xi_3$ characterizes linear polarization along the $\xi$ or $\eta$ axis; the probability of photon polarization along these axes is, respectively, $(1+\xi_3)/2$ or $(1-\xi_3)/2$. Thus $\xi_3=+1$ or $-1$ represents complete polarization in the corresponding direction.

The parameter $\xi_1$ characterizes linear polarization along directions making an angle $\varphi=\pi/4$ or $\varphi=-\pi/4$ with the $\xi$ axis. The respective probabilities of photon polarization along these directions are $(1+\xi_1)/2$ and $(1-\xi_1)/2$; this follows immediately by projecting $\rho_{\alpha\beta}$ on the directions $\mathbf{e}=(1,\pm1)/\sqrt{2}$.

Finally, $\xi_2$ is the degree of circular polarization; by (8.6), the probabilities that the photon has right- or left-circular polarization are $(1+\xi_2)/2$ and $(1-\xi_2)/2$, respectively. Since the two polarizations correspond to helicities $\lambda=\pm1$, it follows that, in the general case, $\xi_2$ is the mean photon helicity. For a pure state with polarization $\mathbf{e}$, note also that
\begin{equation}
\xi_2=i(\mathbf{e}\times\mathbf{e}^*)\mathbf{n}.
\tag{8.10}
\end{equation}

Recall (see II, \S~50) that $\xi_2$ and $\sqrt{\xi_1^2+\xi_3^2}$ are invariant under Lorentz transformations.

We shall also need below the behavior of the Stokes parameters under time reversal. It is readily seen that they are invariant under this transformation. This property clearly does not depend on the nature of the polarization state, so it is enough to verify it for a pure state. In quantum mechanics, time reversal corresponds to replacing the wave function by its complex conjugate (see III, \S~18). For a plane-polarized wave this gives the replacement\footnote{The additional sign change of $\mathbf{e}$ arises because time reversal changes the sign of the electromagnetic vector potential. The scalar potential does not change sign; hence time reversal acts on the 4-vector $e$ as
\begin{equation}
(e_0,\mathbf{e})\to(e_0^*,-\mathbf{e}^*).
\tag{8.11a}
\end{equation}}
\begin{equation}
\mathbf{k}\to-\mathbf{k},
\qquad
\mathbf{e}\to-\mathbf{e}^*.
\tag{8.11}
\end{equation}

\SourcePage{40}{45}
Under this transformation, the symmetric part of the density matrix, $(1/2)(e_\alpha e^*_\beta+e_\beta e^*_\alpha)$, and therefore $\xi_1$ and $\xi_3$, remain unchanged. The invariance of $\xi_2$ under the same transformation follows from (8.10); it is also immediate from the meaning of $\xi_2$ as the mean helicity. Indeed, helicity is the projection of the angular momentum $\mathbf{j}$ on the direction $\mathbf{n}$, namely the product $\mathbf{j}\mathbf{n}$, and time reversal reverses both vectors.

In subsequent calculations we shall require the photon density matrix written in four-dimensional form, that is, as a 4-tensor $\rho_{\mu\nu}$. For a polarized photon described by the 4-vector $e_\mu$, this tensor is naturally defined by
\begin{equation}
\rho_{\mu\nu}=e_\mu e^*_\nu.
\tag{8.12}
\end{equation}
In the three-dimensionally transverse gauge $e=(0,\mathbf{e})$, if one spatial coordinate axis is taken along $\mathbf{n}$, the nonzero components of this 4-tensor coincide with (8.7).

For an unpolarized photon, the three-dimensionally transverse gauge gives a tensor $\rho_{\mu\nu}$ whose components are
\begin{equation}
\rho_{ik}=\frac{1}{2}(\delta_{ik}-n_i n_k),
\qquad
\rho_{0i}=\rho_{i0}=\rho_{00}=0
\tag{8.13}
\end{equation}
(if one axis coincides with $\mathbf{n}$, this reduces to (8.8)). Direct use of $\rho_{\mu\nu}$ in this three-dimensional form is inconvenient, however. We may instead make a gauge transformation; for the density matrix it has the form
\begin{equation}
\rho_{\mu\nu}\to\rho_{\mu\nu}+\chi_\mu k_\nu+\chi_\nu k_\mu,
\tag{8.14}
\end{equation}
where the $\chi_\mu$ are arbitrary functions. Setting
\[
\chi_0=-1/(4\omega),
\qquad
\chi_i=k_i/(4\omega^2),
\]
we obtain, in place of (8.13), the simple four-dimensional expression
\begin{equation}
\rho_{\mu\nu}=-g_{\mu\nu}/2.
\tag{8.15}
\end{equation}

The four-dimensional representation of the density matrix for a partially polarized photon is readily constructed by first rewriting the two-dimensional tensor (8.9) in three-dimensional form:
\[
\begin{split}
\rho_{ik}={}&(1/2)(e_i^{(1)}e_k^{(1)}+e_i^{(2)}e_k^{(2)})
+(\xi_1/2)(e_i^{(1)}e_k^{(2)}+e_i^{(2)}e_k^{(1)})-{}\\
&(i\xi_2/2)(e_i^{(1)}e_k^{(2)}-e_i^{(2)}e_k^{(1)})
+(\xi_3/2)(e_i^{(1)}e_k^{(1)}-e_i^{(2)}e_k^{(2)}),
\end{split}
\]
where $\mathbf{e}^{(1)}$, $\mathbf{e}^{(2)}$ are unit vectors along the $\xi$ and $\eta$ axes. The required generalization is obtained by replacing these 3-vectors with real spacelike unit 4-vectors $e^{(1)}$, $e^{(2)}$ that are mutually orthogonal and orthogonal to the photon 4-momentum $k$:
\begin{equation}
e^{(1)2}=e^{(2)2}=-1,
\qquad
e^{(1)}e^{(2)}=0,
\qquad
e^{(1)}k=e^{(2)}k=0.
\tag{8.16}
\end{equation}

\SourcePage{41}{46}
In the three-dimensionally transverse gauge, $e^{(1)}=(0,\mathbf{e}^{(1)})$, $e^{(2)}=(0,\mathbf{e}^{(2)})$. Thus the four-dimensional photon density matrix is
\begin{equation}
\begin{split}
\rho_{\mu\nu}={}&(1/2)(e_\mu^{(1)}e_\nu^{(1)}+e_\mu^{(2)}e_\nu^{(2)})
+(\xi_1/2)(e_\mu^{(1)}e_\nu^{(2)}+e_\mu^{(2)}e_\nu^{(1)})-{}\\
&(i\xi_2/2)(e_\mu^{(1)}e_\nu^{(2)}-e_\mu^{(2)}e_\nu^{(1)})
+(\xi_3/2)(e_\mu^{(1)}e_\nu^{(1)}-e_\mu^{(2)}e_\nu^{(2)}).
\end{split}
\tag{8.17}
\end{equation}
The most convenient specific choice of the 4-vectors $e^{(1)}$, $e^{(2)}$ depends on the conditions of the problem under consideration.

It must be borne in mind that (8.16) does not fix $e_\mu^{(1)}$ and $e_\mu^{(2)}$ uniquely. If a 4-vector $e_\mu$ satisfies these conditions, then so does every 4-vector $e_\mu+\chi k_\mu$ (because $k^2=0$). This nonuniqueness is connected with the gauge nonuniqueness of the density matrix.

The first term in (8.17) represents the unpolarized state. It could therefore be replaced, in accordance with (8.15), by $-g_{\mu\nu}/2$. Such a replacement is again equivalent to a gauge transformation. In operations with 4-tensors of the form (8.17), expanded in terms of two independent 4-vectors, the following formal device is useful. Writing (8.17) as
\[
\rho_{\mu\nu}=\sum_{a,b=1}^{3}\rho^{(ab)}e_\mu^{(a)}e_\nu^{(b)},
\]
represent the coefficients $\rho^{(ab)}$ as a two-row matrix,
\[
\rho=
\begin{pmatrix}
\rho^{(11)} & \rho^{(12)}\\
\rho^{(21)} & \rho^{(22)}
\end{pmatrix}.
\]
Like any Hermitian two-row matrix, it may be expanded in four independent two-row matrices: the Pauli matrices $\sigma_x$, $\sigma_y$, $\sigma_z$, and the unit matrix~1:
\begin{equation}
\rho=(1/2)(1+\boldsymbol{\xi}\boldsymbol{\sigma}),
\qquad
\boldsymbol{\xi}=(\xi_1,\xi_2,\xi_3),
\tag{8.18}
\end{equation}
as follows directly by comparison with (8.17), using the familiar expressions (18.5) for the Pauli matrices. (Combining the three quantities $\xi_1$, $\xi_2$, $\xi_3$ into a ``vector'' $\boldsymbol{\xi}$ is, of course, purely formal and serves only to make the notation compact.)

\begin{center}
\ProblemHeading
\end{center}

Write the photon density matrix in the representation whose coordinate ``axes'' are the circular unit vectors (8.2).

\SolutionHeading Components of the tensor $\rho'_{\alpha\beta}$ in the new axes ($\alpha$, $\beta=\pm1$) are obtained by projecting the tensor (8.9) onto the unit vectors (8.2):
\[
\rho'_{11}=\rho_{\alpha\beta}e_\alpha^{(+1)*}e_\beta^{(+1)},
\qquad
\rho'_{-1-1}=\rho_{\alpha\beta}e_\alpha^{(-1)*}e_\beta^{(-1)},
\]
\[
\rho'=\frac{1}{2}
\begin{pmatrix}
1+\xi_2 & -\xi_3+i\xi_1\\
-\xi_3-i\xi_1 & 1-\xi_2
\end{pmatrix}.
\]
